
\def\PUDcodigo{EME121}

\def\PUDcodacad{04506.39}

\def\PUDdisciplina{Transferência de Calor}

\def\PUDsemestre{S8}

\def\PUDhoras{80}

%NOVOS ---------------------------------------------------
\def\PUDhorasTeoria{78}
\def\PUDhorasPratica{2}

\def\PUDinterECA{ 
\hyperref[PUD:EME115]{Mecanismos (04506.25)}
    
\hyperref[PUD:EME114]{Termodinâmica (04506.27)}
    
\hyperref[PUD:EME119]{Sistemas Mecânicos (04506.35)}
}

\def\PUDatuaisECA{---}

\def\PUDrecursos{Projetor multimídia; Quadro branco e pincel; Aulas em laboratório.}

%----------------------------------------------------------

\def\PUDcreditos{4}

\def\PUDprerequisito{ \hyperref[PUD:EME114]{EME114} }

\def\PUDpreacad{ \hyperref[PUD:EME114]{Termodinâmica (04506.27)} }

\def\PUDobjetivos{Ser capaz de delinear os fenômenos de transporte pertinentes para qualquer processo ou sistema envolvendo transferência de calor; Ser capaz de usar as informações necessárias para calcular taxas de transferência de calor e/ou temperaturas de materiais; Ser capaz de desenvolver modelos representativos de processos ou sistemas reais e tirar conclusões sobre oprojeto ou o desempenho de processos/sistemas a partir da respectiva análise.}

\def\PUDementa{Revisão da Termodinâmica; Condução; Condução Unidimensional em Regime Estacionário; Condução Bidimensional em Regime Estacionário; Condução Transiente; Convecção; Escoamento Externo; Escoamento Interno; Convecção Natural;
Ebulição e Condensação; Trocadores de Calor; Radiação.}

\def\PUDconteudo{ & Revisão da Termodinâmica: Equações de taxas, Conservação da energia, Análise de problemas de transferência de calor (6 horas-aula). 
\\ 
 & Condução: Equação da taxa da condução, As propriedades térmicas da matéria, A equação da difusão de calor, Condições de contorno e inicial (6 horas-aula). 
\\ 
 & Condução Unidimensional em Regime Estacionário: A parede plana, Sistemas radiais, Condução com geração de enegia térmica, Transferância de calor em superfícies estendindas, A equação do Biocalor (8 horas-aula). 
\\ 
 & Condução Bidimensional em Regime Estacionário: O método da separação de variáveis, O fator de forma da condução e a taxa de condução de calor adimensional, Equações diferenciais finitas, O método gráfico (8 horas-aula). 
\\ 
 & Condução Transiente: O método da capacitância global, Validade do método da capacitância, Análise geral via capacitância global, Efeitos espaciais, A parede plana com convecção, Sistemas radiais com convecção, O sólido semi-finito, Objetos com temperaturas ou fluxos térmicos, Constantes na superfície, Aquecimento periódico, Métodos de diferenças finitas, (8 horas-aula). 
\\ 
 & Convecção: As camadas limite da convecção, Coeficientes convectivos local e médio, Escoamentos laminar e turbulento, As equações de camada-limite, Similaridade na camada limite, Significado físico dos parâmetros adimensionais, Analogias das camadas limite, O coeficientes convectivos (6 horas-aula). 
\\ 
 & Escoamento Externo: O método empírico, A placa plana em ecoamento paralelo, Metodologia para um cálculo de convecção, O cilindro em escoamento cruzado, A esfera, Escoamento externo em matrizes tubulares, Jatos colidentes, Leitos recheados (6 horas-aula). 
\\ 
 & Escoamento Interno: Considerações fluidodinâmicas, Considerações térmicas, O balanço de energia, Escoamento laminar em tubos circulares, Correlações da convecção, Intensificação da transferência de calor, Escoamento interno em microescala, Transferência de massa por convecção (6 horas-aula). 
\\ 
 & Convecção Natural: As equações da convecção natural, Considerações de similaridade, Convecção natural sobre uma superfície vertical, Os efeitos da turbulência, Correlações empíricas, Conecção natural no interior de canais formados entres placas paralelas, Convecção natural e forçada combinadas (8 horas-aula). 
\\ 
 & Ebulição e Condensação: Parâmetros adimensionais na ebulição e na condensação, Modos de ebulição, Ebulição em piscina, Correlações da ebulição em piscina, Ebulição com convecção forçada, Condensação, Condensação em gotas (6 horas-aula). 
\\ 
 & Trocadores de Calor: Tipos de trocadores de calor, O coeficiente global de transferência de calor, Análise de trocadores de calor, Cálculo de projeto e de desempenho de trocadores de calor, Trocadores de calor compactos (6 horas-aula). 
\\ 
 & Radiação: Conceitos fundamentais, Intensidade de radiação, Radiação de corpo negro, Emissão de superfícies reais, Absorção, Reflexão e transmissão em superfícies reais, Lei de Kirchhoff, A superfície cinza, Radiação ambiental(6 horas-aula). }

\def\PUDmetodo{Aulas expositórias que podem ser teóricas e/ou práticas, onde as práticas no laboratório serão marcadas no decorrer da disciplina. Um projeto de trocador de calor, com a construção do mesmo, será cobrado logo início da disciplina.}

\def\PUDavalia{A avaliação será feita com aplicação de provas teóricas e/ou práticas, além da possibilidade de inclusão de trabalhos e seminários no decorrer da disciplina. Uma das avaliações será do projeto com o devido memorial de cálculo e construção experimental.
}

\def\PUDbibbasica{ & \href{www.biblioteca.ifce.edu.br/index.asp?codigo_sophia=61984}{Incropera}, F. P.; Witt, D. P.  Fundamentos de Transferência de Calor e de Massa.  7ª ed.  Rio de Janeiro, RJ: LTC, 2014.  672p.  ISBN: 9788521625049. 
\\ 
 &  \href{www.biblioteca.ifce.edu.br/index.asp?codigo_sophia=62558}{Cengel}, Y. A.; Ghajar, A. J.  Transferência de calor e massa: uma abordagem prática.  4ª ed.  São Paulo, SP: MacGraw-Hill, 2012.  902p.  ISBN: 9788580551273. 
%\\ 
 %&  BEJAN, A.;  Transferência de Calor.  1ª edição.  Editora: EDGAR BLUCHER.  ISBN : 9788521200260. }
\\
 & \href{www.biblioteca.ifce.edu.br/index.asp?codigo_sophia=24058}{Moran}, M. J.; Shapiro, H. N.; Muson, B. R.; DeWitt, D. P.  Introdução à engenharia de sistemas térmicos: termodinâmica, mecânica dos fluidos e transferência de calor.  Rio de Janeiro, RJ: LTC, 2005.  604p.  ISBN: 8521614463. }

\def\PUDbibcomplementar{ & \href{www.biblioteca.ifce.edu.br/index.asp?codigo_sophia=63133}{Fox}, R. W.; MacDolnad, A. T.; Pritchard, P. J.  Introdução à Mecânica dos Fluidos.  8ª ed.  Rio de Janeiro, RJ: LTC, 2014.  871p.  ISBN: 9788521623021. 
\\ 
 &  \href{www.biblioteca.ifce.edu.br/index.asp?codigo_sophia=65080}{Nussenzveig}, H. Moysés.  Curso de física básica 2 : fluidos, oscilações e ondas, calor.  5º ed.  São Paulo, SP : Edgard Blücher, 2014.  375p.  ISBN: 9788521207474.
\\
 &  \href{www.biblioteca.ifce.edu.br/index.asp?codigo_sophia=62609}{Young}, Hugh D.  Física II: termodinâmica e ondas.  14º ed.  São Paulo, SP : Pearson Education do Brasil, 2016.  374p.  ISBN: 9788543005737.
\\
 &  \href{www.biblioteca.ifce.edu.br/index.asp?codigo_sophia=40619}{Livi}, Celso Pohlmann.  Fundamentos de fenômenos de transporte: um texto para cursos básicos.  2º ed.  Rio de Janeiro, RJ: LTC, 2012.  237p.  ISBN: 9788521620570.
\\
 &  \href{www.biblioteca.ifce.edu.br/index.asp?codigo_sophia=65410}{Borgnakke}, Claus.  Fundamentos da termodinâmica.  São Paulo, SP: Edgard Blücher, 2013.  728p.  ISBN: 9788521207924. }
%  Maliska, C. R.;  Transferência de Calor e Mecânica dos Fluidos Computacional, 2ª edição, Editora: LTC.  ISBN : 9788521613961.

\def\PUDautor{Francisco Frederico dos Santos Matos}

\def\PUDdata{2014-04-23}

\def\PUDrevisao{3}

\def\PUDrevisor{Francisco Frederico dos Santos Matos}

\def\PUDdatarevisao{2019-05-15}

\PUDcorpo
